\TieudeT{PHONG TỎA VẬT LÍ}
\subsubsection*{PHẦN I. Thí sinh trả lời câu hỏi từ câu 1 đến câu 18. Mỗi câu hỏi thí sinh chỉ chọn một phương án}
\setcounter{ex}{0}
\begin{ex}%book
Trường hợp nào sau đây quỹ đạo của vật là đường thẳng?
\choice
{Một học sinh đi xe từ nhà đến trường}
{\True Một viên bi sắt được thả rơi}
{Một ôtô chuyển động trên đường}
{Một viên đá được ném theo phương ngang}

\loigiai{
}
\haidong{10} \end{ex}
\begin{ex}%book
Độ lớn độ dịch chuyển phụ thuộc vào
\choice
{\True khoảng cách vị trí đầu và vị trí cuối}
{thời điểm xuất phát}
{khối lượng của vật}
{kích thước của vật}
\loigiai{
}
\haidong{10} \end{ex}
\begin{ex}%book
Một chất điểm chuyển động trên trục $Ox$ từ tọa độ $x_1$ đến tọa độ $x_2$. Độ dịch chuyển của chất điểm này là
\choice
{$d=x_1-x_2$}
{\True $d=x_2-x_1$}
{$d=x_1+x_2$}
{$d=x_1.x_2$}
\loigiai{
}
\haidong{10} \end{ex}
\begin{ex}%book
Đối với vật chuyển động, quãng đường đi được
\choice
{là đại lượng vectơ}
{\True có đơn vị đo là m (mét)}
{có đơn vị là $m/s$ (mét trên giây)}
{có đơn vị là s (giây)}
\loigiai{
}
\haidong{10} \end{ex}
\begin{ex}%book
Độ dịch chuyển và quãng đường đi được bằng nhau
\choice
{\True khi vật chuyển động thẳng, không đổi chiều}
{khi vật chuyển động thẳng, đổi chiều}
{khi vật chuyển động thẳng}
{xảy ra ở mọi trường hợp}

\loigiai{
}
\haidong{10} 
\end{ex}
\begin{ex}%book%[1P1Y1-1][Loai1]
Trường hợp nào dưới đây có thể coi vật chuyển động như một chất điểm?
\choice
{Quyển sách rơi từ trên bàn xuống sàn nhà}
{Chiếc ô tô trong bến xe}
{\True Mặt Trăng trong chuyển động quanh Trái Đất}
{Con cá trong chậu nước}
\loigiai{
Kích thước Mặt Trăng so với khoảng cách giữa Trái Đất và Mặt Trăng là rất nhỏ nên Mặt Trăng trong trường hợp này có thể coi như là một chất điểm.
}
\haidong{10} \end{ex}	
\begin{ex}%book%[1P1Y1-1][Loai1]
Trong trường hợp nào sau đây \textbf{không} thể coi vật chuyển động như một chất điểm?
\choice
{Trái Đất trong chuyển động quanh Mặt Trời}
{Ô tô đi từ Hà Nội đến Quảng Ninh}
{Viên đạn chuyển động trong không khí}
{\True Người ngư dân di chuyển trên chiếc thuyền đánh cá}
\loigiai{
Kích thước người ngư dân so với thuyền đánh cá là đáng kể nên không thể coi người ngư dân trong trường hợp này như là chất điểm.
}
\haidong{10} \end{ex}	
\begin{ex}%book
\impicinpar{Số liệu về độ dịch chuyển và thời gian của chuyển động thẳng của một xe ô tô đồ chơi chạy bằng pin được ghi trong bảng bên. Dựa vào bảng này để xác định từ giây thứ 3 đến 5 là xe chuyển động như thế nào?}
{\begin{tabular}{|c|c|c|c|c|c|c|}
\hline
Độ dịch chuyển (m) & 1 & 3 & 5 & 7 & 7 & 7 \\
\hline
Thời gian (s) & 0 & 1 & 2 & 3 & 4 & 5 \\
\hline
\end{tabular}}

\choice
{Chuyển động nhanh dần}
{Chuyển động chậm dần}
{Chuyển động đều}
{\True Không chuyển động}

\loigiai{
}
\haidong{10} \end{ex}
\begin{ex}%book
\immini[thm]{Một vận động viên chạy từ một siêu thị (A) đến cổng Sân Vận Động (D) theo hai quỹ đạo khác nhau. Hãy so sánh độ dịch chuyển và quãng đường chạy được của người vận động viên trong trường hợp đi từ ABCD.
\choice
{$d>s$}
{$d=s$}
{\True $d<s$}
{$d=s=0$}}
{\includegraphics[width=8.5cm]{images/4_2}}
\loigiai{
}
\haidong{10} \end{ex}
\begin{ex}%book%[1P1Y1-1][Loai1]
Trường hợp nào dưới đây có thể coi vật là chất điểm?
\choice
{Trái Đất trong chuyển động tự quay quanh mình nó}
{Hai hòn bi lúc va chạm với nhau}
{Người nhảy cầu lúc đang rơi xuống nước}
{\True Giọt nước mưa lúc đang rơi}
\loigiai{
Giọt nước mưa có kích thước rất nhỏ so với đường đi của nó nên có thể coi là chất điểm.
}
\haidong{10} \end{ex}
\begin{ex}%book %[1P1Y1-1][Loai1] 
Trong các trường hợp sau đây, trường hợp nào có thể xem vật như một chất điểm?
\choice
{Trái Đất đang chuyển động tự quay quanh mình nó}
{Viên đạn đang chuyển động trong nòng súng}
{Tàu hoả đứng trong sân ga}
{\True Trái Đất chuyển động trên quỹ đạo quanh Mặt Trời}
\loigiai{
Trong trường hợp Trái Đất chuyển động trên quỹ đạo quanh Mặt Trời thì có thể coi Trái Đất là chất điểm vì kích thước của Trái Đất là rất nhỏ so với quỹ đạo chuyển động quanh Mặt Trời.
}
\haidong{10} \end{ex}
\begin{ex}%book %[1P1B1-1][Loai1] 
Để xác định vị trí một ô tô du lịch thám hiểm giữa sa mạc người ta sử dụng
\choice
{khoảng cách theo đường chim bay thẳng đến một thành phố ven sa mạc}
{\True kinh độ và vĩ độ địa lí của xe}
{quãng đường đi của xe từ lúc xuất phát ở một thành phố ven sa mạc}
{khoảng thời gian đã đi của xe từ lúc xuất phát đó}
\loigiai{
\begin{itemize}
\item Xác định vị trí trên mặt đất, trên mặt biển người ta dùng kinh độ và vĩ độ địa lí là phổ biến và thuận tiện nhất. 
\item Ngày nay người ta dùng các vệ tinh địa tĩnh trong hệ thống GPS định vị toàn cầu cho nhanh và tiện lợi.
\end{itemize}
}
\haidong{10} 
\end{ex}
\begin{ex}%book
Một người lái xe ô tô đi thẳng $6\ km$ theo hướng Tây, sau đó rẽ trái đi thẳng theo hướng Nam $4\ km$ rồi quay sang hướng Đông đi $3\ km$. Quãng đường đi được của ô tô là
\choice
{$16\ km$}
{\True $13\ km$}
{$10\ km$}
{$14\ km$}
\haidong{12}
\loigiai{
}
\end{ex}
\begin{ex}%book
Bạn Việt đi xe đạp đến trường theo đường thẳng như hình vẽ. Khi đến quầy báo, bạn sực nhớ đã quên một cuốn sách ở nhà nên quay về nhà lấy sách rồi lại đạp xe đến trường. Độ dịch chuyển khi bạn quay từ quầy báo về nhà và độ dịch chuyển khi bạn từ quầy báo đến trường lần lượt là
\begin{center}
\begin{tikzpicture}[draw=Mapcolor, very thick,scale=0.6,>=stealth']
\draw[->](0,0)--++(0:12);
\draw[fill=white](0,0)circle(1.5pt)node[above]{O (Nhà)};
\draw[fill=white](4,0)circle(1.5pt)node[above]{B (Quầy báo)}node[below]{400m};
\draw[fill=white](10,0)circle(1.5pt)node[above]{T (Trường học)}node[below]{1000m};
\end{tikzpicture}
\end{center}
\choice
{$400\ m$ và $600\ m$}
{$400\ m$ và $1000\ m$}
{$-400\ m$ và $1000\ m$}
{\True $-400\ m$ và $600\ m$}
\haidong{12}
\loigiai{
\begin{itemize}
\item Độ dịch chuyển khi bạn Việt từ quầy báo về nhà: $s_2=\overline{BO}=-\overline{OB}=-x_B=-400\ m$.
\item Độ dịch chuyển khi bạn Việt từ quầy báo đến trường: $s_3=\overline{BT}=\overline{OT}-\overline{OB}=x_T-x_B=1000-400=600\ m$.
\end{itemize}
}
\end{ex}

\begin{ex}%book
Một ôtô đi $12\ km$ theo hướng tây và sau đó đi $16\ km$ về hướng nam. Quãng đường đi được và độ dịch chuyển của ô tô là
\choice
{$28\ km$ và $28\ km$}
{\True $28\ km$ và $20\ km$}
{$20\ km$ và $28\ km$}
{$20\ km$ và $20\ km$}
\haidong{12}
\loigiai{
}
\end{ex}
\begin{ex}%book
Một người lái ô tô đi thẳng $6\ km$ theo hướng Tây, sau đó rẽ trái đi thẳng theo hướng Nam $8\ km$. Quãng đường $s$ đi được và độ dịch chuyển $d$ của ô tô lần lượt là
\choice
{\True $14\ km;\ 10\ km$}
{$10\ km;\ 14\ km$}
{$10\ km;\ 10\ km$}
{$14\ km;\ 14\ km$}
\haidong{12}
\loigiai{
}
\end{ex}
\begin{ex}%Tailieuchuan
Một vật chuyển động thẳng trên một đoạn đường. Nếu chọn $10$ giờ làm mốc tính thời gian và thời điểm $11$ giờ lúc xuất phát thì thời điểm bắt đầu đo thời gian là
\choice
{$30$ phút}
{\True $1$ giờ}
{$2$ giờ}
{$4$ giờ}
\loigiai{

}
\end{ex}
\begin{ex}%Tailieuchuan
Một người điều khiển thuyền đi được $5,6\ km$ theo hướng Bắc trên mặt hồ phẳng lặng trong thời gian $1,0\ h$. Sau đó, anh ta quay thuyền đi về phía tây $3,4\ km$ trong $30$ phút. Độ dịch chuyển tổng hợp của thuyền có độ lớn
\choice
{\True $6,6\ km$ và hướng lệch về phía Tây $31^\circ$ so với hướng Bắc}
{$13,2\ km$ và hướng lệch về phía Tây $31^\circ$ so với hướng Bắc}
{$13,2\ km$ và hướng lệch về phía Đông $42^\circ$ so với hướng Bắc}
{$6,6\ km$ và hướng lệch về phía Đông $42^\circ$ so với hướng Bắc}
\loigiai{

}
\end{ex}











\subsubsection*{PHẦN 2. Thí sinh trả lời từ câu 1 đến câu 4. Trong mỗi ý a), b), c), d) ở mỗi câu, thí sinh chọn đúng hoặc sai.}
\setcounter{ex}{0}

\begin{ex}%book
Trong hình bên người đi xe máy $(1)$, người đi bộ $(2)$, người đi ôtô $(3)$ đều khởi hành từ siêu thị $A$ để đi đến bưu điện $B$.
\immini{\choiceTF[t]
{\True Độ lớn độ dịch chuyển của ba người là bằng nhau}
{Quãng đường đi được của người đi bộ là lớn nhất trong ba người}
{\True Độ lớn của độ dịch chuyển và quãng đường đi được bằng nhau nếu người đó đi theo đường thẳng từ A đến B}
{Quãng đường đi được luôn nhỏ hơn độ dịch chuyển}
}{
\begin{tikzpicture}
\node[] at (0,0){\includegraphics[width=7cm]{images/4-6}};
\node[] at (0,-4.5){\textit{Sơ đồ mô tả quãng đường đi được của người}};
\node[] at (0,-5){\textit{đi xe máy, người đi bộ và người đi ô tô}};
\end{tikzpicture}}
\loigiai{
\begin{itemchoice}
\itemch 
\itemch
\itemch 
\itemch
\end{itemchoice}
}
\end{ex}
\begin{ex}%Tailieuchuan
\immini[thm]{Độ lớn độ dịch chuyển mô tả ở hình vẽ trong hệ tọa độ địa lí. Biết chuyển động là chuyển động thẳng. Tỉ xích $1\ cm$ ứng với $100\ m$.}{\begin{tikzpicture}[scale=1,thick,>=stealth']
\begin{scope}[shift={(0,3.5)},rotate=45,scale=0.4,draw=Mapcolor]
\tkzDefPoints{0/0/o, 1/0/d, -1/0/t, 0/1/b, 0/-1/n}
\tkzDefShiftPoint[o](45:0.2cm){db}
\tkzDefShiftPoint[o](135:0.2cm){tb}
\tkzDefShiftPoint[o](225:0.2cm){tn}
\tkzDefShiftPoint[o](315:0.2cm){dn}
\fill[red](o)--(d)--(db)--cycle;
\fill[red](o)--(b)--(tb)--cycle;
\fill[red](o)--(t)--(tn)--cycle;
\fill[red](o)--(n)--(dn)--cycle;
\fill[black](o)--(d)--(dn)--cycle;
\fill[black](o)--(b)--(db)--cycle;
\fill[black](o)--(t)--(tb)--cycle;
\fill[black](o)--(n)--(tn)--cycle;
\end{scope}
\begin{scope}[shift={(0,3.5)},rotate=0,scale=0.6]
\tkzDefPoints{0/0/o, 1/0/d, -1/0/t, 0/1/b, 0/-1/n}
\tkzDefShiftPoint[o](45:0.2cm){db}
\tkzDefShiftPoint[o](135:0.2cm){tb}
\tkzDefShiftPoint[o](225:0.2cm){tn}
\tkzDefShiftPoint[o](315:0.2cm){dn}
\fill[red](o)--(d)--(db)--cycle;
\fill[red](o)--(b)--(tb)--cycle;
\fill[red](o)--(t)--(tn)--cycle;
\fill[red](o)--(n)--(dn)--cycle;
\fill[black](o)--(d)--(dn)--cycle;
\fill[black](o)--(b)--(db)--cycle;
\fill[black](o)--(t)--(tb)--cycle;
\fill[black](o)--(n)--(tn)--cycle;
\tkzLabelPoint[right](d){Đ}
\tkzLabelPoint[left](t){T}
\tkzLabelPoint[below](n){N}
\tkzLabelPoint[above](b){B}
\end{scope}

\path
(-5.5,0.5)coordinate(O)
;
\draw[->,Mapcolor] (-5.5,0.5) -- (-6.75,0.5);
\draw[->,Mapcolor] (-5.5,0.5) -- (-2.5,0.5);
\draw[->,Mapcolor] (-5.5,0) -- (-5.5,2.5);
\draw[->,Mapcolor] (-5.5,0.5) -- (-4,2);
\node at (-5.5,1.5) {};
\node at (-4.5,0.5) {};
\node at (-3.5,0.5) {};
\node at (-4.75,1.25) {};
\draw[Mapcolor] (-5.6,1.5) -- (-5.4,1.5);
\draw[Mapcolor] (-4.5,0.6) -- (-4.5,0.4);
\draw[Mapcolor] (-3.5,0.6) -- (-3.5,0.4);
\draw[Mapcolor] (-4.8,1.3) -- (-4.7,1.2);
\draw[thin,Mapcolor] (-5.3,0.7) .. controls (-5.11,0.76) and (-5.05,0.59) .. (-5.15,0.5);
\node at (-5.5,2.8) {$d_1$};
\node at (-3.7,2.2) {$d_2$};
\node at (-2.2,0.5) {$d_3$};
\node at (-7.1,0.5) {$d_4$};
\node[below left] at (-5.5,0.5) {$0$};
\node at (-4.7,0.8) {$45^\circ $};

\end{tikzpicture}}
\choiceTF[t]
{\True Độ lớn độ dịch chuyển $d_{4}=100 \ m$} 
{\True Quãng đường đi được $s_{4}$ bằng độ lớn độ dịch chuyển $d_{4}$} 
{Độ dịch chuyển $d_1=100 \ m$ theo hướng Bắc} 
{\True Độ dịch chuyển $d_2=200 \ m$ hướng theo góc $45^\circ$ từ Đông sang Bắc} 
\loigiai{ 
\begin{itemchoice}
\itemch
\itemch
\itemch
\itemch
\end{itemchoice}
}
\end{ex}
\begin{ex}%Tailieuchuan
Một người đi thang máy từ tầng $G$ xuống tầng hầm cách tầng $G$ $5\ m$, rồi lên tới tầng cao nhất của toà nhà cách tầng G $50\ m$.
\choiceTF[t]
{Khi đi từ tầng $G $xuống tầng hầm thì quãng đường đi được là $-5\ m$} 
{\True Khi đi từ tầng $G$ xuống tầng hầm độ lớn độ dịch chuyển là $5\ m$} 
{\True Khi đi từ tầng hầm lên tầng cao nhất quãng đường đi được là $55\ m$} 
{\True Khi đi từ tầng hầm lên tầng cao nhất độ dịch chuyển có độ lớn là $55\ m$} 
\loigiai{ 
\begin{itemchoice}
\itemch
\itemch
\itemch
\itemch
\end{itemchoice}
}
\end{ex}
\begin{ex}%Tailieuchuan
Một người đi bộ $5,0 \ km$ trên một con đường thẳng theo hướng Bắc rồi quay đầu lại và đi $12\ km$ theo hướng Nam theo sơ đồ như hình.
\immini{\choiceTF[t]
{\True Vectơ độ dịch chuyển đổi chiều trong quá trình chuyển động} 
{Tổng quãng đường đã đi là $12\ km$} 
{\True Quãng đường đi được và độ lớn của độ dịch chuyển trong trường hợp trên là khác nhau} 
{\True Độ lớn độ dịch chuyển là $7\ km$ theo hướng Nam} }{\begin{tikzpicture}[scale=1,thick,>=stealth',draw=Mapcolor]
\path
(0,0)coordinate(O)
;
\node [fill=Mapcolor, draw, thick, shape=circle, minimum width=.25cm, minimum height=.25cm, anchor=center] at (O) {};
\draw[->,Mapcolor] (0,0) -- (0,1.5);
\draw[->,Mapcolor] (0.5,1.5) -- (0.5,-2);
\node at (-0.6,0.75) {$5\ km$};
\node at (1.25,-0.25) {$12\ km$};
\end{tikzpicture}}
\loigiai{ 
\begin{itemchoice}
\itemch
\itemch
\itemch
\itemch
\end{itemchoice}
}
\end{ex}


\subsubsection*{PHẦN 3. Thí sinh trả lời từ câu 1 đến câu 6.}
\setcounter{ex}{0}
\textit{Sử dụng thông tin sau cho Câu 1 và Câu 2:} Một người đi xe máy từ nhà đến bến xe bus cách nhà $6\ km$ về phía Đông. Đến bến xe, người đó lên xe bus đi tiếp $20\ km$ về phía Bắc.
\begin{ex}%Tailieuchuan
Quãng đường đi được trong cả chuyến đi là bao nhiêu $km$?
\shortans[oly]{26}
\loigiai{
}
\end{ex}
\begin{ex}%Tailieuchuan
Độ dịch chuyển tổng hợp của người đó có độ lớn là bao nhiêu km (làm tròn kết quả đến chữ số hàng phần mười)?
\shortans[oly]{20,9}
\loigiai{
}
\end{ex}
\textit{Sử dụng thông tin sau cho Câu 3 và Câu 4:} Bạn A đi xe đạp từ nhà qua trạm xăng, tới siêu thị mua đồ rồi quay về nhà cất đồ, sau đó đi xe đến trường.
\begin{center}
\includegraphics[width=12cm]{images/4-7}
\end{center} 
\begin{ex}%Tailieuchuan
Quãng đường đi được từ trạm xăng đến siêu thị là bao nhiêu mét?
\shortans[oly]{400}
\loigiai{
}
\end{ex}
\begin{ex}%Tailieuchuan
Độ dịch chuyển của cả chuyến đi có độ lớn là bao nhiêu mét?
\shortans[oly]{1200}
\end{ex}
\textit{Sử dụng thông tin sau cho Câu 5 và Câu 6:} Một thang máy mang một người từ tầng trệt xuống tầng hầm sâu $5\ m$ rồi đi lên tầng $3$. Biết chiều cao tầng trệt và các lầu là $4\ m$. Chọn gốc tọa độ tại mặt đất.
\begin{ex}%Tailieuchuan
Quãng đường chuyển động khi người này lên tới lầu 3 là bao nhiêu mét?
\shortans[oly]{18}
\loigiai{
}
\end{ex}
\begin{ex}%Tailieuchuan
Độ dịch chuyển từ khi thang máy đi từ tầng hầm đến khi dừng tại lầu 3 là bao nhiêu mét?
\shortans[oly]{13}
\end{ex}
\newpage 
\TieudeT{PHONG TỎA VẬT LÍ}
\subsubsection*{PHẦN I. Thí sinh trả lời câu hỏi từ câu 1 đến câu 18. Mỗi câu hỏi thí sinh chỉ chọn một phương án}
\setcounter{ex}{0}
\begin{ex}%book
Phát biểu nào sau đây đúng khi nói về quỹ đạo của vật?
\choice
{Quỹ đạo là một đường thẳng mà trên đó chất điểm chuyển động}
{Một đường vạch sẵn trong không gian trên đó chất điểm chuyển động gọi là quỹ đạo}
{\True Quỹ đạo là một đường nối những vị trí liên tiếp của vật theo thời gian trong quá trình chuyển động}
{Một đường cong mà trên đó chất điểm chuyển động gọi là quỹ đạo}
\loigiai{
}
\haidong{10} \end{ex}
\begin{ex}%book
Độ dịch chuyển là đại lượng
\choice
{\True vectơ}
{vô hướng}
{luôn âm}
{luôn dương}

\loigiai{
}
\haidong{10} \end{ex}
\begin{ex}%book
Độ dịch chuyển của một vật là
\choice
{quãng đường vật đi được trong khoảng thời gian đó}
{khoảng thời gian vật đi từ điểm đầu đển điểm cuối}
{tốc độ chuyển động nhanh hay chậm của vật}
{\True một vectơ nối vị trí điểm đầu đến mút là điểm cuối của chuyển động và độ dài tỉ lệ với độ lớn của độ dịch chuyển}

\loigiai{
}
\haidong{10} \end{ex}	
\begin{ex}%book
Quãng đường đi được là một đại lượng luôn có giá trị
\choice
{âm}
{dương, âm hoặc bằng không}
{dương hoặc âm}
{\True không âm}
\loigiai{
}
\haidong{10} \end{ex}	
\begin{ex}%book
Phát biểu nào sau đây là đúng khi nói về độ dịch chuyển và quãng đường đi được của một vật?
\choice
{Độ dịch chuyển và quãng đường đi được đều là đại lượng vectơ}
{\True Độ dịch chuyển là đại lượng vectơ còn quãng đường đi được là đại lượng vô hướng}
{Độ dịch chuyển và quãng đường đi được đều là đại lượng vô hướng}
{Độ dịch chuyển và quãng đường đi được đều là đại lượng không âm}

\loigiai{
}
\haidong{10} \end{ex}	
\begin{ex}%book
Phát biểu nào sau đây là đúng khi so sánh quãng đường và độ dịch chuyển? 
\choice
{\True Quãng đường luôn luôn lớn hơn hoặc bằng độ lớn của độ dịch chuyển}
{Quãng đường luôn luôn nhỏ hơn hoặc bằng độ lớn của độ dịch chuyển}
{Quãng đường luôn luôn bằng độ lớn của độ dịch chuyển}
{Quãng đường luôn luôn nhỏ hơn độ lớn của độ dịch chuyển}

\loigiai{
}
\haidong{10} \end{ex}	

\begin{ex}%book %[1P1B1-1][Loai1] 
Để xác định vị trí của một tàu biển giữa đại dương, người ta sử dụng
\choice
{độ sâu của nước biển tại vị trí của tàu}
{\True kinh độ và vĩ độ địa lí của tàu}
{khoảng cách từ tàu đến bến cảng gần nhất}
{khoảng cách từ tàu đến một bến cảng nhất định}
\loigiai{
Để xác định vị trí của một tàu biển giữa đại dương, người ta dùng kinh độ và vĩ độ của tàu.
}
\haidong{10} \end{ex}	
\begin{ex}%book
Một người đi bộ $5\ km$ trên một con đường thẳng theo hướng Bắc rồi quay đầu lại và đi $12\ km$ theo hướng Nam. Độ dịch chuyển của người đó là
\choice
{$17\ km$, hướng Nam}
{$5\ km$, hướng Bắc}
{\True $7\ km$, hướng Nam}
{$7\ km$ hướng Bắc}
\haidong{12}
\loigiai{
}
\end{ex}	
\begin{ex}%book
Một con nhện bò dọc theo các cạnh của một chiếc bàn hình chữ nhật. Biết hai cạnh bàn có chiều dài lần lượt là $0,8\ m$ và $1,2\ m$. Độ dịch chuyển của con nhện khi nó xuất phát từ một đỉnh hình chữ nhật và đi được quãng đường $2,0\ m$ là
\choice
{\True $1,4\ m$}
{$1,5\ m$}
{$1,6\ m$}
{$1,7\ m$}
\haidong{12}
\loigiai{
}
\end{ex}	
\begin{ex}%book
\immini[thm]{Một vật bắt đầu chuyển động từ điểm $O$ về điểm $B$, sau đó chuyển động đến điểm $A$ như hình vẽ. Quãng đường và độ dịch chuyển của vật tương ứng bằng}{\begin{tikzpicture}[draw=Mapcolor]
\tkzDefPoints{0/0/o, 2.5/0/x, -1.5/0/x'}
\tkzDrawSegments[->,Mapcolor](x',x)
\foreach  \x  in  {-0.5,0.5,1}
{\draw[] (\x,-0.1)  --  (\x,0.1);}
\foreach  \x  in  {-1,0,1.5}
{\draw[fill=white] (\x,0) circle(1pt);}
\node[above] at (0,0){$O$};
\node[above] at (-1,0){$B$};
\node[above] at (1.5,0){$A$};
\node[below] at (-1,0){$-2$};
\node[below] at (-0.5,0){$-1$};
\node[below] at (0,0){$0$};
\node[below] at (0.5,0){$1$};
\node[below] at (1,0){$2$};
\node[below] at (1.5,0){$3$};
\tkzLabelPoint[above](x){$x\ (m)$}
\end{tikzpicture}
}
\choice
{\True $7\ m;\ 3\ m$}
{$7\ m;-3\ m$}
{$8\ m;\ 3\ m$}
{$8\ m;-2\ m$}
\haidong{12}
\loigiai{
}
\end{ex}
\begin{ex}%Tailieuchuan
Một người chỉ cho một người khách du lịch như sau: "Ông hãy đi dọc theo phố này đến một bờ hồ lớn. Đứng tại đó, nhìn sang bên kia hồ theo hướng Tây Bắc, ông sẽ thấy tòa nhà của khách sạn S". Người chỉ đường đã xác định vị trí của khách sạn S theo cách nào?
\choice
{Cách dùng đường đi và vật làm mốc}
{Cách dùng hệ tọa độ địa lí}
{\True Dùng cả hai cách A và B}
{Không dùng cả hai cách A và B}
\loigiai{

}
\end{ex}
\begin{ex}%Tailieuchuan
Một người đi thang máy từ tầng G xuống tầng hầm cách tầng G $5\ m$. Quãng đường và độ dịch chuyển của người đó khi đi từ tầng G xuống tầng hầm là
\choice
{Quãng đường $s = 50\ m;\ d = 5\ m$ (xuống dưới)}
{Quãng đường $s = 5\ m;\ d = 5\ m$ (lên trên)}
{Quãng đường $s = 45\ m;\ d = -5\ m$ (xuống dưới)}
{\True Quãng đường $s = 5\ m;\ d = 5\ m$ (xuống dưới)}
\loigiai{

}
\end{ex}


\textit{Sử dụng dữ kiện sau cho Câu 13 đến Câu 18:} Một người di chuyển ở bốn địa điểm: xe buýt đậu bên đường, cây xanh, trường học, công viên.
\begin{center}
\includegraphics[width=12cm]{images/2025_04_24_df450a76ef00368ca939g-29}
\end{center}
\begin{ex}%Tailieuchuan
Quãng đường đi được của người đó di chuyển từ chỗ xe buýt đến trường rồi đến công viên là
\choice
{$80\ m$}
{$100\ m$}
{\True $180\ m$}
{$128\ m$}
\loigiai{

}
\end{ex}
\begin{ex}%Tailieuchuan
Độ lớn độ dịch chuyển khi người đó di chuyển từ chỗ xe buýt đến trường rồi đến công viên là
\choice
{$80\ m$}
{$100\ m$}
{$180\ m$}
{\True $128\ m$}
\loigiai{

}
\end{ex}
\begin{ex}%Tailieuchuan
Quãng đường đi được của người đó di chuyển từ chỗ xe buýt đến trường rồi đến công viên và cuối cùng đến chỗ cây xanh là
\choice
{$80\ m$}
{$100\ m$}
{$180\ m$}
{\True $260\ m$}
\loigiai{

}
\end{ex}
\begin{ex}%Tailieuchuan
Độ lớn độ dịch chuyển khi người đó di chuyển từ chỗ xe buýt đến trường rồi đến công viên và cuối cùng đến chỗ cây xanh là
\choice
{$80\ m$}
{\True $100\ m$}
{$180\ m$}
{$128\ m$}
\loigiai{

}
\end{ex}
\begin{ex}%Tailieuchuan
Quãng đường đi được của người đó di chuyển từ chỗ xe buýt đến trường rồi đến công viên rồi đến chỗ cây xanh và cuối cùng đến chỗ xe buýt là
\choice
{$80\ m$}
{\True $360\ m$}
{$180\ m$}
{$128\ m$}
\loigiai{}
\end{ex}
\begin{ex}%Tailieuchuan
Độ lớn độ dịch chuyển khi người đó di chuyển từ chỗ xe buýt đến trường rồi đến công viên rồi đến chỗ cây xanh và cuối cùng đến chỗ xe buýt là
\choice
{$80\ m$}
{$100\ m$}
{\True $0\ m$}
{$128\ m$}
\loigiai{

}
\end{ex}	
\subsubsection*{PHẦN 2. Thí sinh trả lời từ câu 1 đến câu 4. Trong mỗi ý a), b), c), d) ở mỗi câu, thí sinh chọn đúng hoặc sai.}
\setcounter{ex}{0}
\begin{ex}%Tailieuchuan
Xét một vật chuyển động.
\choiceTF[t]
{\True Quãng đường là đại lương vô hướng, không âm, chỉ được đặc trưng bởi độ lớn} 
{\True Quãng đường đi được không thể hiện được chiều chuyển động của vật} 
{\True Độ dịch chuyển có thể nhận giá trị âm, dương hoặc bằng không} 
{Khi vật chuyển động thẳng, không đổi chiều thì độ lớn của độ dịch chuyển và quãng đường đi được khác nhau} 
\loigiai{ 
\begin{itemchoice}
\itemch
\itemch
\itemch
\itemch
\end{itemchoice}
}
\end{ex}
\begin{ex}%book
\immini[thm]{Hai người đạp xe đạp từ  $A$ đến $C$, người thứ nhất đi theo đường từ $A$ đến $B$, rồi từ $B$ đến $C$, người thứ hai đi thẳng từ $A$ đến $C$ như hình vẽ.}{\begin{tikzpicture}[scale=0.6,thick,>=stealth',draw=Mapcolor]
\tkzDefPoints{0/0/a, 0/3/b, 4/3/c}
\tkzDrawPolygon[Mapcolor](a,b,c)
\tkzLabelSegment[above,sloped](a,b){$3\ km$}
\tkzLabelSegment[above](b,c){$4\ km$}
\tkzMarkRightAngle[size=0.3](a,b,c)
\tkzLabelPoint[below](a){$A$}
\tkzLabelPoint[above](b){$B$}
\tkzLabelPoint[above](c){$C$}
\end{tikzpicture}}
\choiceTF[t]
{\True Người thứ nhất đi được quãng đường $7\ km$, còn độ dịch chuyển là $5\ km$}
{\True Người thứ hai có quãng đường và độ dịch chuyển đều bằng $5\ km$}
{Độ dịch chuyển của người thứ nhất là $7\ km$ theo hướng Đông - Bắc}
{Góc giữa độ dịch chuyển của người thứ nhất và hướng Đông là $53^\circ$}
\loigiai{
\begin{itemchoice}
\itemch $AB = 3\ km$, $BC = 4\ km$ nên quãng đường người thứ nhất: $s_1 = 3 + 4 = 7\ km$
\itemch Do tam giác vuông tại $B$ nên $AC = \sqrt{AB^2 + BC^2} = \sqrt{9 + 16} = 5\ km$
$\Rightarrow$ Độ dịch chuyển của cả hai người đều là $5\ km$
\itemch Độ dịch chuyển người thứ nhất không phải $7\ km$, nên phát biểu 3 sai
\itemch $\cos \angle{BAC} = \dfrac{3}{5} \Rightarrow \angle{BAC} \approx 53^\circ$ là góc giữa $\vec{AC}$ và phương Bắc, không phải phương Đông.
\end{itemchoice}
}
\end{ex}	
\begin{ex}%Tailieuchuan
Em của An chơi trò chơi tìm kho báu ở ngoài vườn với các bạn của mình. Em của An giấu kho báu của mình là một chiếc vòng nhựa vào trong một chiếc giày rồi viết mật thư tìm kho báu như sau: Bắt đầu từ gốc cây ổi, đi $10$ bước về phía Bắc, sau đó đi $4$ bước về phía Tây, $15$ bước về phía Nam, $5$ bước về phía Đông và $5$ bước về phía Bắc là tới chỗ giấu kho báu.
\choiceTF[t]
{Không thể tìm ra kho báu theo cách mô tả này} 
{Quãng đường phải đi (theo bước) đế tìm ra kho báu là 34 bước} 
{Kho báu được giấu ở vị trí cách cây ổi $5$ bước theo hướng nam} 
{\True Độ dịch chuyển (theo bước) để tìm ra kho báu là 1 bước theo hướng đông} 
\loigiai{ 
\begin{itemchoice}
\itemch
\itemch
\itemch
\itemch
\end{itemchoice}
}
\end{ex}
\begin{ex}%Tailieuchuan
Xét hai xe chuyển động như hình. Chọn gốc toạ độ tại vị trí O trên hình vẽ, chiều dương hướng từ trái sang phải. Tại một thời điểm nào đó, xe màu xanh và xe màu cam đang lần lượt ở vị trí có toạ độ $x_{A}=10\ km$ và $x_{B}=20 \ km$. Sau khoảng thời gian $\Delta t_1$, xe màu xanh đến được vị trí $x_{B}$ và xe màu cam đến được vị trí $x_{A}$.
\begin{center} 
\includegraphics[width=10cm]{images/2025_04_24_df450a76ef00368ca939g-16(1)}
\end{center} 
\choiceTF[t]
{\True Hai xe chuyển động ngược chiều nhau} 
{\True Quãng đường đi được của hai xe là $10\ km$} 
{Độ lớn độ dịch chuyển bằng quãng đường đi được của hai xe} 
{Độ dịch chuyển của xe B có giá trị là $10\ km$} 
\loigiai{ 
\begin{itemchoice}
\itemch
\itemch
\itemch
\itemch
\end{itemchoice}
}
\end{ex}

\subsubsection*{PHẦN 3. Thí sinh trả lời từ câu 1 đến câu 6.}
\setcounter{ex}{0}
\immini{\textit{Sử dụng thông tin sau cho Câu 1 và Câu 2:} Một người đi bộ với tốc độ không đổi dọc theo nửa đường tròn có bán kính $5,0\ m$, từ A đến $B$ như hình với thời gian đi là $6,0\ s$.}{\begin{tikzpicture}[scale=1,thick,>=stealth',draw=Mapcolor]
\path
(0,0)coordinate(O)
(4.5,-0.5)coordinate(A)
;

\draw[<->,Mapcolor] (-0.5,0) -- (0.5,0);
\draw[<->,Mapcolor] (0,-0.5) -- (0,0.5);
\draw[<-,Mapcolor] (A) arc (0:180:1 and 1);
\node at (2.5,-1) {$A$};
\node at (4.5,-1) {$B$};
\node at (0,1) {$B$};
\node at (0,-1) {$N$};
\node at (1,0) {$\text{Đ}$};
\node at (-1,0) {$T$};
\end{tikzpicture}}
\begin{ex}%Tailieuchuan
Độ lớn độ dịch chuyển là bao nhiêu mét?
\shortans[oly]{10}
\loigiai{
}
\end{ex}
\begin{ex}%Tailieuchuan
Quãng đường đã đi là bao nhiêu mét (làm tròn kết quả đến chữ số hàng phần mười)?
\shortans[oly]{15,7}
\loigiai{
}
\end{ex}
\textit{Sử dụng thông tin sau cho Câu 3 và Câu 4:} \immini{Hai người đi xe đạp từ $A$ đến $C$, người thứ nhất đi theo đường từ $A$ đến $B$, rồi từ $B$ đến $C$; người thứ hai đi thẳng từ $A$ đến $C$ như hình vẽ. Cả hai đều về đích cùng một lúc.
}{\begin{tikzpicture}[scale=0.6,thick,>=stealth',draw=Mapcolor]
\tkzDefPoints{0/0/a, 0/3/b, 4/3/c}
\tkzDrawPolygon[Mapcolor](a,b,c)
\tkzLabelSegment[above,sloped](a,b){$4\ km$}
\tkzLabelSegment[above](b,c){$4\ km$}
\tkzMarkRightAngle[size=0.3](a,b,c)
\tkzLabelPoint[below](a){$A$}
\tkzLabelPoint[above](b){$B$}
\tkzLabelPoint[above](c){$C$}
\end{tikzpicture}}
\begin{ex}%Tailieuchuan
Quãng đường đi được của người thứ nhất là bao nhiêu mét?
\shortans[oly]{8}
\loigiai{
}
\end{ex}
\begin{ex}%Tailieuchuan
Độ lớn độ dịch chuyển của người thứ hai là bao nhiêu mét (làm tròn kết quả đến chữ số hàng phần mười)?
\shortans[oly]{5,7}
\loigiai{}
\end{ex}
\textit{Sử dụng thông tin sau cho Câu 5 và Câu 6:} Một người lái ô tô đi thẳng $6\ km$ theo hướng Tây, sau đó rẽ trái đi thẳng theo hướng Nam $4\ km$ rồi quay sang hướng Đông đi $3\ km$.
\begin{ex}%Tailieuchuan
Quãng đường đi được của ô tô là bao nhiêu km?
\shortans[oly]{13}
\loigiai{}
\end{ex}
\begin{ex}%Tailieuchuan
Độ lớn độ dịch chuyển của ô tô là bao nhiêu km?
\shortans[oly]{5}
\loigiai{}
\end{ex}